% GENERATED FILE. Edit canonical lessons, then run npm run generate:books.
% canonical-source-hash: fnv1a64:40488d5dffeca637
% canonical-lessons: BN-C517-laphdori, BN-C517-dhaoya-kora, BN-C517-bedminton, BN-C517-lukochuri, BN-C517-marbel

\chapter{Games Outdoors}
\label{ch:bn-games-outdoors}
\hlchaptermodality{517}

\begin{chapteropening}
\textbf{By the end of this chapter:} \emph{I can talk about games outdoors, skipping and jumping over things.}

Complete the last lesson of chapter 517: I can talk about games outdoors, skipping and jumping over things.
\end{chapteropening}

\section[\texorpdfstring{\bn{লাফদড়ি} (lāphdôṛi)}{লাফদড়ি (lāphdôṛi)}]{\bn{লাফদড়ি} (lāphdôṛi) — a skipping rope; skipping}

\label{lesson:BN-C517-laphdori}



\begin{quote}
Before the new one: say the Bengali for absorbed, then the Bengali for amateur.
\end{quote}

\subsection*{You'll want to know: \bn{লাফদড়ি}}
\textbf{\bn{লাফদড়ি}} — \emph{lāphdôṛi} — \textquotedblleft{}a skipping rope; skipping\textquotedblright{}.

\begin{quote}
\bn{মেয়েরা} \bn{উঠোনে} \bn{লাফদড়ি} \bn{খেলছিল।}
\end{quote}

\begin{quote}
\emph{Meyerā uṭhone lāphdôṛi khelchhilô.}
\end{quote}

\begin{quote}
The girls were skipping in the courtyard.
\end{quote}

\begin{cousinweb}
\textbf{\bn{লাফদড়ি}} (\emph{lāphdôṛi}) joins \bn{লাফ}, \textquotedblleft{}a leap\textquotedblright{}, from a reconstructed \emph{lappa} (CDIAL 10939), to \bn{দড়ি}, \textquotedblleft{}a rope\textquotedblright{}.
\end{cousinweb}

\begin{grammarlens}[title={What you've built}]
One more word for describing what you lived and felt.
\end{grammarlens}

\subsection*{Your turn}
\begin{itemize}
\raggedright
  \item \emph{Say it:} \emph{lāphdôṛi}
  \item \emph{Say it:} \emph{lāphdôṛi}, once more
  \item \emph{Say it:} \emph{Meyerā uṭhone lāphdôṛi khelchhilô.} — The girls were skipping in the courtyard.
  \item \emph{Recall:} say the Bengali for absorbed, then the Bengali for amateur, then say \emph{lāphdôṛi} again
\end{itemize}

\begin{tcolorbox}[breakable,colback=teal!4,colframe=teal!35!black,title={Before you move on}]
What does \bn{লাফদড়ি} mean? (\textquotedblleft{}A skipping rope; skipping\textquotedblright{}.) Say it once more.
\end{tcolorbox}
\section[\texorpdfstring{\bn{ধাওয়া} \bn{করা} (dhāoyā kôrā)}{ধাওয়া করা (dhāoyā kôrā)}]{\bn{ধাওয়া} \bn{করা} (dhāoyā kôrā) — to chase; to run after}

\label{lesson:BN-C517-dhaoya-kora}



\begin{quote}
Before the new one: say the Bengali for amateur, then the Bengali for a skipping rope.
\end{quote}

\subsection*{You'll want to know: \bn{ধাওয়া} \bn{করা}}
\textbf{\bn{ধাওয়া} \bn{করা}} — \emph{dhāoyā kôrā} — \textquotedblleft{}to chase; to run after\textquotedblright{}.

\begin{quote}
\bn{কুকুরটা} \bn{বিড়ালটাকে} \bn{ধাওয়া} \bn{করল।}
\end{quote}

\begin{quote}
\emph{Kukurṭā biṛālṭāke dhāoyā kôrlô.}
\end{quote}

\begin{quote}
The dog chased the cat.
\end{quote}

Past (he/she …): \bn{তিনি} \bn{চোরটাকে} \bn{ধাওয়া} \bn{করলেন}, \emph{tini chorṭāke dhāoyā kôrlen}.

\begin{cousinweb}
\textbf{\bn{ধাওয়া} \bn{করা}} (\emph{dhāoyā kôrā}) joins \bn{ধাওয়া}, \textquotedblleft{}a run, a chase\textquotedblright{}, to \bn{করা}, \textquotedblleft{}to do\textquotedblright{}; CDIAL lists Bengali \emph{dhāoyā}, \textquotedblleft{}to run\textquotedblright{}, under Sanskrit \emph{dhāvati}, \textquotedblleft{}runs\textquotedblright{} (CDIAL 6802).
\end{cousinweb}

\begin{grammarlens}[title={What you've built}]
One more word for describing what you lived and felt.
\end{grammarlens}

\subsection*{Your turn}
\begin{itemize}
\raggedright
  \item \emph{Say it:} \emph{dhāoyā kôrā}
  \item \emph{Say it:} \emph{dhāoyā kôrā}, once more
  \item \emph{Say it:} \emph{Kukurṭā biṛālṭāke dhāoyā kôrlô.} — The dog chased the cat.
  \item \emph{Recall:} say the Bengali for amateur, then the Bengali for a skipping rope, then say \emph{dhāoyā kôrā} again
\end{itemize}

\begin{tcolorbox}[breakable,colback=teal!4,colframe=teal!35!black,title={Before you move on}]
What does \bn{ধাওয়া} \bn{করা} mean? (\textquotedblleft{}To chase; to run after\textquotedblright{}.) Say it once more.
\end{tcolorbox}
\section[\texorpdfstring{\bn{ব্যাডমিন্টন} (bæḍminṭôn)}{ব্যাডমিন্টন (bæḍminṭôn)}]{\bn{ব্যাডমিন্টন} (bæḍminṭôn) — badminton}

\label{lesson:BN-C517-bedminton}



\begin{quote}
Before the new one: say the Bengali for a skipping rope, then the Bengali for to run after.
\end{quote}

\subsection*{You'll want to know: \bn{ব্যাডমিন্টন}}
\textbf{\bn{ব্যাডমিন্টন}} — \emph{bæḍminṭôn} — \textquotedblleft{}badminton\textquotedblright{}.

\begin{quote}
\bn{শীতের} \bn{সন্ধ্যায়} \bn{আমরা} \bn{পাড়ায়} \bn{ব্যাডমিন্টন} \bn{খেলি।}
\end{quote}

\begin{quote}
\emph{Shīter sôndhyāy āmrā pāṛāy bæḍminṭôn kheli.}
\end{quote}

\begin{quote}
On winter evenings we play badminton in the neighbourhood.
\end{quote}

\begin{cousinweb}
\textbf{\bn{ব্যাডমিন্টন}} (\emph{bæḍminṭôn}) is borrowed from English \emph{badminton}.
\end{cousinweb}

\begin{grammarlens}[title={What you've built}]
One more word for describing what you lived and felt.
\end{grammarlens}

\subsection*{Your turn}
\begin{itemize}
\raggedright
  \item \emph{Say it:} \emph{bæḍminṭôn}
  \item \emph{Say it:} \emph{bæḍminṭôn}, once more
  \item \emph{Say it:} \emph{Shīter sôndhyāy āmrā pāṛāy bæḍminṭôn kheli.} — On winter evenings we play badminton in the neighbourhood.
  \item \emph{Recall:} say the Bengali for a skipping rope, then the Bengali for to run after, then say \emph{bæḍminṭôn} again
\end{itemize}

\begin{tcolorbox}[breakable,colback=teal!4,colframe=teal!35!black,title={Before you move on}]
What does \bn{ব্যাডমিন্টন} mean? (\textquotedblleft{}Badminton\textquotedblright{}.) Say it once more.
\end{tcolorbox}
\section[\texorpdfstring{\bn{লুকোচুরি} (lukochuri)}{লুকোচুরি (lukochuri)}]{\bn{লুকোচুরি} (lukochuri) — hide-and-seek}

\label{lesson:BN-C517-lukochuri}



\begin{quote}
Before the new one: say the Bengali for to run after, then the Bengali for badminton.
\end{quote}

\subsection*{You'll want to know: \bn{লুকোচুরি}}
\textbf{\bn{লুকোচুরি}} — \emph{lukochuri} — \textquotedblleft{}hide-and-seek\textquotedblright{}.

\begin{quote}
\bn{বিকেলে} \bn{আমরা} \bn{উঠোনে} \bn{লুকোচুরি} \bn{খেললাম।}
\end{quote}

\begin{quote}
\emph{Bikele āmrā uṭhone lukochuri khellām.}
\end{quote}

\begin{quote}
In the afternoon we played hide-and-seek in the courtyard.
\end{quote}

\begin{cousinweb}
\textbf{\bn{লুকোচুরি}} (\emph{lukochuri}) joins \bn{লুকো}-, the stem of \bn{লুকানো}, \textquotedblleft{}to hide\textquotedblright{} (CDIAL 11083), to \bn{চুরি}, \textquotedblleft{}theft\textquotedblright{} (CDIAL 4937): hiding and stealing away.
\end{cousinweb}

\begin{grammarlens}[title={What you've built}]
One more word for describing what you lived and felt.
\end{grammarlens}

\subsection*{Your turn}
\begin{itemize}
\raggedright
  \item \emph{Say it:} \emph{lukochuri}
  \item \emph{Say it:} \emph{lukochuri}, once more
  \item \emph{Say it:} \emph{Bikele āmrā uṭhone lukochuri khellām.} — In the afternoon we played hide-and-seek in the courtyard.
  \item \emph{Recall:} say the Bengali for to run after, then the Bengali for badminton, then say \emph{lukochuri} again
\end{itemize}

\begin{tcolorbox}[breakable,colback=teal!4,colframe=teal!35!black,title={Before you move on}]
What does \bn{লুকোচুরি} mean? (\textquotedblleft{}Hide-and-seek\textquotedblright{}.) Say it once more.
\end{tcolorbox}
\section[\texorpdfstring{\bn{মার্বেল} (mārbel)}{মার্বেল (mārbel)}]{\bn{মার্বেল} (mārbel) — a marble (for games)}

\label{lesson:BN-C517-marbel}



\begin{quote}
Before the new one: say the Bengali for badminton, then the Bengali for hide-and-seek.
\end{quote}

\subsection*{You'll want to know: \bn{মার্বেল}}
\textbf{\bn{মার্বেল}} — \emph{mārbel} — \textquotedblleft{}a marble (for games)\textquotedblright{}.

\begin{quote}
\bn{ছেলেটা} \bn{পকেটে} \bn{পাঁচটা} \bn{মার্বেল} \bn{রাখল।}
\end{quote}

\begin{quote}
\emph{Chheleṭā pôkeṭe pā̃chṭā mārbel rākhlô.}
\end{quote}

\begin{quote}
The boy put five marbles in his pocket.
\end{quote}

\begin{cousinweb}
\textbf{\bn{মার্বেল}} (\emph{mārbel}) is borrowed from English \emph{marble}.
\end{cousinweb}

\begin{grammarlens}[title={What you've built}]
One more word for describing what you lived and felt.
\end{grammarlens}

\subsection*{Your turn}
\begin{itemize}
\raggedright
  \item \emph{Say it:} \emph{mārbel}
  \item \emph{Say it:} \emph{mārbel}, once more
  \item \emph{Say it:} \emph{Chheleṭā pôkeṭe pā̃chṭā mārbel rākhlô.} — The boy put five marbles in his pocket.
  \item \emph{Recall:} say the Bengali for badminton, then the Bengali for hide-and-seek, then say \emph{mārbel} again
\end{itemize}

\begin{tcolorbox}[breakable,colback=teal!4,colframe=teal!35!black,title={Before you move on}]
What does \bn{মার্বেল} mean? (\textquotedblleft{}A marble (for games)\textquotedblright{}.) Say it once more.
\end{tcolorbox}
